\input format.tex

\usepackage{graphicx}
\graphicspath{{cores/}}

\begin{document}

\vspace*{6mm}
%% 各章节
\setlength{\arrayrulewidth}{1pt}
\fontsize{9.3pt}{11pt}\selectfont
\color{gray2}

{\noindent\bf\sanhao 一、肠道菌群概况}

\vspace*{6mm}
\begin{spacing}{1.25}
\begin{LRaside}[.8]{\fontsize{11pt}{11pt}\selectfont 1.1 肠道菌群多样性}
\noindent
\includegraphics[width=\linewidth]{diversity.pdf}
\asidebreak %
肠道菌群是一个复杂而精巧的生态系统，不同菌群之间相互竞争，共同维持着肠道生态系统的稳定。\\
肠道菌群多样性：通过检测肠道内上千种细菌的种类和含量，从而全面反映菌群生态平衡的指标。多样性越高，菌株之间就越容易维持动态平衡，菌群生态系统就越稳定，越不容易被外界的因素（如不规律的饮食、药物滥用等）所破坏。
\end{LRaside}
\end{spacing}

\vspace*{6mm}

\begin{spacing}{1.25}
\begin{LRaside2}{\bf 您本次肠道菌群多样性检测结果}
\begin{center}
\setlength{\unitlength}{1cm}
\begin{picture}(20,2.6)(0,-0.8)
\put(3.35132258064516,0.55){\includegraphics[scale=1]{straw_l4.pdf}}
\put(0.5,-0.2){\includegraphics[scale=1]{duoyangxing_bar.pdf}}
\end{picture}
\indent\fontsize{9pt}{11pt}\selectfont\bf {综合您本次有益菌检测结果，您肠道菌群多样性检测结果为
{\fontsize{13pt}{14pt}\selectfont\color{level2} 2.30}，需要关注
}
\includegraphics[scale=1]{underline.pdf}
\end{center}
\vspace{1mm}
\asidebreak
\begin{tabular}{p{1.5cm}p{13cm}}
\begin{minipage}{1.4cm}{\vspace{1mm}\includegraphics[scale=0.8]{result.pdf}} \end{minipage}
& \parbox[c]{\hsize}{\vskip7pt {您的肠道菌群多样性偏高，说明您的肠道菌群组成十分丰富，菌群失调风险低。您的肠道菌群不易受到外界因素的干扰，但仍需减少不良作息、刺激性食物、药物滥用等的影响，这样有助于保持良好的肠道菌群多样性，增强机体新陈代谢和免疫能力。} \vskip7pt} \\
\end{tabular}

\end{LRaside2}
\end{spacing}

\vspace*{-2mm}
\hspace*{12.5cm}
\raisebox{-0.4ex}{\includegraphics{drop_l4.pdf}}
\fontsize{9pt}{11pt}\selectfont {本次检测结果}

\end{document}

